%% Commands for TeXCount
%TC:macro \cite [option:text,text]
%TC:macro \citep [option:text,text]
%TC:macro \citet [option:text,text]
%TC:envir table 0 1
%TC:envir table* 0 1
%TC:envir tabular [ignore] word
%TC:envir displaymath 0 word
%TC:envir math 0 word
%TC:envir comment 0 0
%%
%%
%% For submission and review:
%% \documentclass[acmsmall,screen,review,anonymous]{acmart}.
\documentclass[acmsmall,screen,nonacm]{acmart}

% \setcopyright{rightsretained}
% \acmDOI{10.1145/3674642}
% \acmYear{2024}
% \acmJournal{PACMPL}
% \acmVolume{8}
% \acmNumber{ICFP}
% \acmArticle{253}
% \acmMonth{8}
% \acmSubmissionID{icfp24main-p58-p}
% \received{2024-02-28}
% \received[accepted]{2024-06-18}

\citestyle{acmnumeric}

\usepackage{mathpartir}
\usepackage{wrapfig}
\usepackage{scalerel}
\usepackage{minted}
\usepackage{mathtools}
\usepackage{fontawesome} % lock symbol
\usepackage{tikz}
% \usepackage{quiver}
\usepackage{stmaryrd}
\usepackage{dsfont}
\usepackage{pifont}

\newcommand{\under}[2]{\underset{\raisebox{2mm}{$\scriptstyle #2$}}{#1}}
\newcommand{\inh}[1]{\under{#1}{\uparrow}}
\newcommand{\syn}[1]{\under{#1}{\downarrow}}
\DeclareMathOperator{\llet}{\text{let}}
\DeclareMathOperator{\lin}{\text{in}}
\DeclareMathOperator{\case}{\text{case}}
\DeclareMathOperator{\bbox}{\text{box}}
\DeclareMathOperator{\unbox}{\text{unbox}}
\DeclareMathOperator{\mmod}{\text{mod}}
\newcommand{\unmod}[3]{\text{let}_{#1} \text{ mod}_{#2}({#3}) \leftarrow}
\newcommand{\at}{\mathbin{\mbox{\small\ensuremath{@}}}}
\DeclarePairedDelimiter\ann{\langle}{\rangle}
\DeclareMathOperator{\return}{\text{return}}
\newcommand{\unit}{\mathds{1}}
\DeclareMathOperator{\inl}{\text{inl}}
\DeclareMathOperator{\inr}{\text{inr}}
\DeclareMathOperator{\suc}{\text{succ}}
\DeclareMathOperator{\primrec}{\text{primrec}}
\newcommand{\nat}{\mathbb{N}}
\newcommand{\tee}{\vdash}
\newcommand{\teev}{\vdash^\mathrm{v}}
\newcommand{\teec}{\vdash^\mathrm{c}}
\newcommand{\hole}{\,?}
\newcommand*{\Ev}{E_\mathrm{v}}
\newcommand*{\Ep}{K}
\DeclareMathOperator{\red}{\rightsquigarrow}
\DeclarePairedDelimiter\floor{\lfloor}{\rfloor}
\DeclarePairedDelimiter\ceil{\lceil}{\rceil}
\DeclareMathOperator{\lock}{\tikz{\node[scale=1, line width=0pt, inner sep=0pt] at (0,0) {\faLock\,};}}
\definecolor{gainsboro}{rgb}{0.86, 0.86, 0.86}
% \newcommand{\gray}[1]{\colorbox{gainsboro}{$#1$}} % {\strut$#1$}}
\newcommand{\reducedstrut}{\vrule width 0pt height .9\ht\strutbox depth .9\dp\strutbox\relax}
\newcommand{\gray}[1]{%
  \begingroup
  \setlength{\fboxsep}{0pt}%
  \colorbox{gainsboro}{\reducedstrut$#1$\/}%
  \endgroup
}

\newcommand{\clashes}{\,\text{\Lightning}\,}
\newcommand{\notclashes}{\,{\text{\Lightning}\mkern-8.5mu\backslash}\,}
\DeclareMathOperator{\joinctx}{\mathbin{\fatsemi}}
\DeclareMathOperator{\splitctx}{\curlyveeuparrow}

\newcommand{\synth}{\Longrightarrow}
\newcommand{\checks}{\Longleftarrow}
\newcommand{\lolli}{\multimap}
\newcommand{\withtype}{\mathbin{\&}}
\newcommand{\susp}{\mathbf{susp}}
\newcommand{\force}{\mathbf{force}}
\newcommand{\down}{\mathbf{down}}
\newcommand{\matches}{\mathbf{match}}
\newcommand{\DW}{\Delta_{\mathrm{W}}}
\newcommand{\one}{\mathbf{1}}
\newcommand{\sem}[1]{\llbracket #1 \rrbracket}

\begin{document}

\title{Internship Report: Modal Type Systems}

\author{Izzy Gupta}
\affiliation{%
  \institution{University of Edinburgh}
  %\city{Edinburgh}
  \country{United Kingdom}
}

% \author{Anton Lorenzen}
% \orcid{0000-0003-3538-9688}
% \affiliation{%
  % \institution{University of Edinburgh}
  % %\city{Edinburgh}
  % \country{United Kingdom}
% }
% \email{anton.lorenzen@ed.ac.uk}

% \author{Sam Lindley}
% \orcid{0000-0002-1360-4714}
% \affiliation{%
  % \institution{University of Edinburgh}
  % %\city{Edinburgh}
  % \country{United Kingdom}
% }
% \email{Sam.Lindley@ed.ac.uk}

%% By default, the full list of authors will be used in the page
%% headers. Often, this list is too long, and will overlap
%% other information printed in the page headers. This command allows
%% the author to define a more concise list
%% of authors' names for this purpose.
% \renewcommand{\shortauthors}{}

\begin{abstract}
  Internship Report
\end{abstract}

%% The code below is generated by the tool at http://dl.acm.org/ccs.cfm.
% \begin{CCSXML}
% <ccs2012>
   % <concept>
       % <concept_id>10003752.10003790.10003793</concept_id>
       % <concept_desc>Theory of computation~Modal and temporal logics</concept_desc>
       % <concept_significance>500</concept_significance>
       % </concept>
   % <concept>
       % <concept_id>10003752.10003790.10003801</concept_id>
       % <concept_desc>Theory of computation~Linear logic</concept_desc>
       % <concept_significance>300</concept_significance>
       % </concept>
 % </ccs2012>
% \end{CCSXML}

% \ccsdesc[500]{Theory of computation~Modal and temporal logics}
% \ccsdesc[300]{Theory of computation~Linear logic}

%% Keywords. The author(s) should pick words that accurately describe
%% the work being presented. Separate the keywords with commas.
% \keywords{Modal Types, Mode Qualifiers}

\maketitle

\section{Calculi}

\subsection{Graded Modal Type Theory}

A fine-grained call-by-value calculus in the style of \cite{Abel:unified,Choudhury:graded}.

\paragraph{Syntax}

\begin{align*}
  A &::= \,^q A \rightarrow B \mid \boxempty^q A \mid A \times B \mid A + B \mid \unit \\
  v &::= x \mid \lambda x.\, e \mid \bbox_q\, v \mid (v_1, v_2) \mid () \mid \inl\, v \mid \inr\, v \\
  e &::= v \mid e\, v
      \mid \llet x \leftarrow e_1 \lin e_2 \\
    &\mid \llet_{q_1} \bbox_{q_2} x = v \lin e \\
    &\mid \llet_q (x, y) = v \lin e
      \mid \case_q\, v\, \{x.e_1;\; y.e_2\}
\end{align*}

\paragraph{Value typing rules} \hfill \framebox{$\Gamma \tee v : A$}

\begin{mathpar}
  \inferrule*[right=var]{
    \phantom{.}
  }{
    0\cdot\Gamma, x :^1 A, 0\cdot\Gamma' \tee x : A
  } \and
  \inferrule*[right=sub]{
    q \leq q' \\
    \Gamma, x :^q A, \Gamma' \tee v : B
  }{
    \Gamma, x :^{q'} A, \Gamma' \tee v : B
  } \and
  \inferrule*[right=lam]{
    \Gamma, x :^q A \tee e : B
  }{
    \Gamma \tee \lambda x.\, e : \,^q A \rightarrow B
  } \and
  \inferrule*[right=unit]{
    \phantom{.}
  }{
    \Gamma \tee () : \unit
  } \and
  \inferrule*[right=pair]{
    \Gamma_1 \tee v_1 : A \\
    \Gamma_2 \tee v_2 : B
  }{
    \Gamma_1 + \Gamma_2 \tee (v_1, v_2) : A \times B
  } \and
  \inferrule*[right=inl]{
    \Gamma \tee v : A
  }{
    \Gamma \tee \inl\, v : A + B
  } \and
  \inferrule*[right=inr]{
    \Gamma \tee v : B
  }{
    \Gamma \tee \inr\, v : A + B
  } \and
  \inferrule*[right=$\boxempty$I]{
    \Gamma \tee v : A
  }{
    q\cdot\Gamma \tee \bbox_q\, v : \boxempty^q A
  }
\end{mathpar}

\paragraph{Expression typing rules} \hfill \framebox{$\Gamma \tee e : B$}

\begin{mathpar}
  \inferrule*[right=app]{
    \Gamma_1 \tee e : \,^q A \rightarrow B \\
    \Gamma_2 \tee v : A
  }{
    \Gamma_1 + q\cdot\Gamma_2 \tee e\, v : B
  } \and
  \inferrule*[right=bind]{
    \Gamma_1 \tee e_1 : A \\
    \Gamma_2, x :^1 A \tee e_2 : B
  }{
    \Gamma_1 + \Gamma_2 \tee \llet x \leftarrow e_1 \lin e_2 : B
  } \and
  \inferrule*[right=$\boxempty$E]{
    \Gamma_1 \tee v : \boxempty^{q_2} A \\
    \Gamma_2, x :^{q_1 \cdot q_2} A \tee e : B
  }{
    q_1\cdot\Gamma_1 + \Gamma_2 \tee \llet_{q_1} \bbox_{q_2} x = v \lin e : B
  } \and
  \inferrule*[right=split]{
    \Gamma_1 \tee v : A_1 \times A_2 \\
    \Gamma_2, x :^q A_1, y :^q A_2 \tee e : B
  }{
    q\cdot\Gamma_1 + \Gamma_2 \tee \llet_q (x, y) = v \lin e : B
  } \and
  \inferrule*[right=case]{
    \Gamma_1 \tee v : A + B \\
    \Gamma_2, x :^q A \tee e_1 : C \\
    \Gamma_2, y :^q B \tee e_2 : C
  }{
    q\cdot\Gamma_1 + \Gamma_2 \tee \case_q\, v\, \{x.e_1;\; y.e_2\} : C
  }
\end{mathpar}

\paragraph{Semantics} \hfill \framebox{$e \red e'$}

\begin{align*}
  (\lambda x.\, e)\, v &\red e[v/x] \\
  \llet x \leftarrow v \lin e &\red e[v/x] \\
  \llet_{q_1} \bbox_{q_2} x = \bbox_{q_2}\, v \lin e &\red e[v/x] \\
  \llet_q (x, y) = (v_1, v_2) \lin e &\red e[v_1/x, v_2/y] \\
  \case_q\, (\inl\, v)\, \{x.e_1;\; y.e_2\} &\red e_1[v/x] \\
  \case_q\, (\inr\, v)\, \{x.e_1;\; y.e_2\} &\red e_2[v/y]
\end{align*}

\paragraph{Evaluation contexts}\mbox{}\\
\noindent
\begin{minipage}[t]{0.55\textwidth}
\begin{align*}
  E &::= [\cdot] \mid E\, v \mid \llet x \leftarrow E \lin e
\end{align*}
\end{minipage}%
\begin{minipage}[t]{0.45\textwidth}
\begin{mathpar}
  \inferrule*[right=ctx]{
    e \red e'
  }{
    E[e] \red E[e']
  }
\end{mathpar}
\end{minipage}

\paragraph{Parameterisation}

Graded modal type theory is usually parameterised by a pre-ordered semiring $(S, 0, 1, +, \cdot, \leq)$ of grades.
Aside from the usual semiring axioms, we require that $+$ and $\cdot$ are monotone with respect to $\leq$.

For local and global modes, we take the three-element semiring $\{0, 1, \mathbb{G}\}$ with
$0 \leq 1 \leq \mathbb{G}$, $+$ as maximum, and $\mathbb{G} \cdot \mathbb{G} = \mathbb{G}$.

\paragraph{Derivations}

We derive terms for the following types:
%
\begin{align*}
  T &= \lambda x.\, \llet_1 \bbox_\mathbb{G} y = x \lin y
    : \,^1 (\boxempty^\mathbb{G} A) \rightarrow A \\
  4 &= \lambda x.\, \llet_1 \bbox_\mathbb{G} y = x \lin \bbox_\mathbb{G}\, (\bbox_\mathbb{G}\, y)
    : \,^1 (\boxempty^\mathbb{G} A) \rightarrow \boxempty^\mathbb{G}\boxempty^\mathbb{G} A \\
  N &= \bbox_q\, () : \boxempty^q \unit \\
  X &= \lambda p.\, \llet_1 (x, y) = p \lin
    \llet_1 \bbox_q f = x \lin
    \llet_1 \bbox_q g = y \lin {} \\
    &\phantom{= \lambda p.\,}
    \bbox_q\, (\lambda a.\,
      \llet b \leftarrow f\, a \lin g\, b) \\
    &\phantom{=}\; : \,^1 (\boxempty^q (\,^1 A \rightarrow B)
      \times \boxempty^q (\,^1 B \rightarrow C)) \\
    &\phantom{=}\; \rightarrow \boxempty^q (\,^1 A \rightarrow C) \\
  M^{\times} &= \lambda p.\, \llet_1 (x, y) = p \lin
    \llet_1 \bbox_q a = x \lin \\
    &\phantom{= \lambda p.\,}
    \llet_1 \bbox_q b = y \lin \bbox_q\, (a, b) \\
    &\phantom{=}\; : \,^1 (\boxempty^q A \times \boxempty^q B) \rightarrow \boxempty^q (A \times B) \\
  M^{+} &= \lambda s.\, \case_1\, s\,
    \{x.\llet_1 \bbox_q a = x \lin \bbox_q\, (\inl\, a);\; \\
    &\phantom{= \lambda s.\, \case\, s\, \{}
    y.\llet_1 \bbox_q b = y \lin \bbox_q\, (\inr\, b)\} \\
    &\phantom{=}\; : \,^1 (\boxempty^q A + \boxempty^q B) \rightarrow \boxempty^q (A + B)
\end{align*}

\subsection{Mode Qualifiers}

\paragraph{Syntax}

\begin{align*}
  A, B &::= A \at \mu_1 \rightarrow B \at \mu_2 \mid \boxempty^\nu A \mid A \times B \mid A + B \mid \unit \\
  e &::= x \mid \lambda x.\, e \mid e_1\, e_2
      \mid \bbox_\nu\, e \mid \unbox_\nu\, e
      \mid (e_1, e_2) \mid () \mid \inl\, e \mid \inr\, e \\
    &\mid \llet x = e_1 \lin e_2
      \mid \llet_\mu (x, y) = e \lin e
      \mid \case_\mu\, e\, \{x.e_1;\; y.e_2\} \\
  \Gamma &::= \cdot \mid \Gamma, x : A \at \mu, \lock_\mu
\end{align*}

We define $\mathrm{locks}(\Gamma)$ as the composition of all locks in $\Gamma$:
\begin{align*}
  \mathrm{locks}(\cdot) &= \mathrm{id} \\
  \mathrm{locks}(\Gamma, x :_\mu A) &= \mathrm{locks}(\Gamma) \\
  \mathrm{locks}(\Gamma, \lock_\mu) &= \mu \circ \mathrm{locks}(\Gamma)
\end{align*}

\paragraph{Typing rules} \hfill \framebox{$\Gamma \tee e : A \at \mu$}

\begin{mathpar}
  \inferrule*[right=var]{
    \mu_1 \leq \mathrm{locks}(\Gamma') \circ \mu_2
  }{
    \Gamma, x : A \at \mu_1, \Gamma' \tee x : A \at \mu_2
  } \and
  \inferrule*[right=lam]{
    \Gamma, \lock_\mu, x : A \at \mu_1 \tee e : B \at \mu_2
  }{
    \Gamma \tee \lambda x.\, e : (A \at \mu_1 \rightarrow B \at \mu_2) \at \mu
  } \and
  \inferrule*[right=app]{
    \Gamma \tee e_1 : (A \at \mu_1 \rightarrow B \at \mu_2) \at \_ \\\\
    \Gamma \tee e_2 : A \at \mu_1
  }{
    \Gamma \tee e_1\, e_2 : B \at \mu_2
  } \and
  \inferrule*[right=$\boxempty$I]{
    \Gamma \tee e : A \at \nu \circ \mu
  }{
    \Gamma \tee \bbox_\nu\, e : \boxempty^\nu A \at \mu
  } \and
  \inferrule*[right=$\boxempty$E]{
    \Gamma \tee e : \boxempty^\nu A \at \mu
  }{
    \Gamma \tee \unbox_\nu\, e : A \at \nu \circ \mu
  } \and
  \inferrule*[right=let]{
    \Gamma \tee e_1 : A \at \mu_1 \\
    \Gamma, x : A \at \mu_1 \tee e_2 : B \at \mu_2
  }{
    \Gamma \tee \llet x = e_1 \lin e_2 : B \at \mu_2
  } \and
  \inferrule*[right=unit]{
    \phantom{.}
  }{
    \Gamma \tee () : \unit \at \mu
  } \and
  \inferrule*[right=pair]{
    \Gamma \tee e_1 : A \at \mu \\
    \Gamma \tee e_2 : B \at \mu
  }{
    \Gamma \tee (e_1, e_2) : A \times B \at \mu
  } \and
  \inferrule*[right=split]{
    \Gamma \tee e_1 : A \times B \at \mu_1 \\\\
    \Gamma, x : A \at \mu_1, y : B \at \mu_1 \tee e_2 : C \at \mu_2
  }{
    \Gamma \tee \llet_{\mu_1} (x, y) = e_1 \lin e_2 : C \at \mu_2
  } \and
  \inferrule*[right=inl]{
    \Gamma \tee e : A \at \mu
  }{
    \Gamma \tee \inl\, e : A + B \at \mu
  } \and
  \inferrule*[right=inr]{
    \Gamma \tee e : B \at \mu
  }{
    \Gamma \tee \inr\, e : A + B \at \mu
  } \and
  \inferrule*[right=case]{
    \Gamma \tee e : A + B \at \mu_1 \\\\
    \Gamma, x : A \at \mu_1 \tee e_1 : C \at \mu_2 \\\\
    \Gamma, y : B \at \mu_1 \tee e_2 : C \at \mu_2
  }{
    \Gamma \tee \case_{\mu_1}\, e\, \{x.e_1;\; y.e_2\} : C \at \mu_2
  }
\end{mathpar}

\paragraph{Call-by-value semantics} \hfill \framebox{$e \red e'$}

\begin{align*}
  v &::= x \mid \lambda x.\, e \mid \bbox_\nu\, v \mid (v_1, v_2) \mid () \mid \inl\, v \mid \inr\, v
\end{align*}

\begin{align*}
  (\lambda x.\, e)\, v &\red e[v/x] \\
  \llet x = v \lin e &\red e[v/x] \\
  \unbox_\nu\, (\bbox_\nu\, v) &\red v \\
  \llet_\mu (x, y) = (v_1, v_2) \lin e &\red e[v_1/x, v_2/y] \\
  \case_\mu\, (\inl\, v)\, \{x.e_1;\; y.e_2\} &\red e_1[v/x] \\
  \case_\mu\, (\inr\, v)\, \{x.e_1;\; y.e_2\} &\red e_2[v/y]
\end{align*}

\paragraph{Evaluation contexts}\mbox{}\\
\noindent
\begin{minipage}[t]{0.6\textwidth}
\begin{align*}
  E &::= [\cdot] \mid \bbox_\nu\, E \mid (E, e) \mid (v, E) \mid \inl\, E \mid \inr\, E \\
    &\mid E\, e_2 \mid v_1\, E \mid \unbox_\nu\, E \\
    &\mid \llet x = E \lin e \mid \llet_\mu (x, y) = E \lin e \\
    &\mid \case_\mu\, E\, \{x.e_1;\; y.e_2\}
\end{align*}
\end{minipage}%
\begin{minipage}[t]{0.4\textwidth}
\begin{mathpar}
  \inferrule*[right=ctx]{
    e \red e'
  }{
    E[e] \red E[e']
  }
\end{mathpar}
\end{minipage}

\paragraph{Parameterisation}

Mode qualifier systems are usually parameterised by a set of axes $\mathcal{A}$.
Each axis $((A, \leq) \in \mathcal{A})$ is a total order and contains a distinguished element $1_A$.
The distinguished element $1_A$ is either the minimum element of $A$ (in which we call this axis \textit{monadic})
or the maximum element of $A$ (in which case we call this axis \textit{comonadic}).
Composition is given by maximum on monadic axes and by minimum on comonadic axes.
A mode $\mu$ is a product of axes which is partially ordered by the pointwise extension of the orders on the axes.

For the local and global modes, we take a single comonadic axis $\{1, \mathbb{G}\}$ and $\mathbb{G} \leq 1$.
For this parameterisation, it is also possible to define the lock as a left-division operator:
\begin{align*}
  \cdot,\lock_\mu &= \cdot \\
  \Gamma, x :_{\mu_1} A, \lock_{\mu_2} &= \Gamma, \lock_{\mu_2}, x :_{\mu_1} A & (\text{if } \mu_1 \leq \mu_2) \\
  \Gamma, x :_{\mu_1} A, \lock_{\mu_2} &= \Gamma, \lock_{\mu_2} & (\text{otherwise})
\end{align*}

\paragraph{Derivations}

We derive terms for the following types:
%
\begin{align*}
  T &= \lambda x.\, \unbox_\mathbb{G}\, x
    : ((\boxempty^\mathbb{G} A) \at 1 \rightarrow A \at 1) \at 1 \\
  4 &= \lambda x.\, \bbox_\mathbb{G}\, (\bbox_\mathbb{G}\, (\unbox_\mathbb{G}\, x))
    : ((\boxempty^\mathbb{G} A) \at 1 \rightarrow (\boxempty^\mathbb{G}\boxempty^\mathbb{G} A) \at 1) \at 1 \\
  N &= \bbox_\nu\, () : \boxempty^\nu \unit \at 1 \\
  X &= \lambda p.\, \llet_{1} (f, g) = p \lin \llet f' = \unbox_\nu\, f \lin \llet g' = \unbox_\nu\, g \lin {} \\
    &\phantom{= \lambda p.\,}
    \bbox_\nu\, (\lambda a.\, g'\, (f'\, a)) \\
    &\phantom{=}\; : ((\boxempty^\nu (A \at 1 \rightarrow B \at 1))
      \times (\boxempty^\nu (B \at 1 \rightarrow C \at 1)) \at 1 \\
    &\phantom{=}\; \rightarrow (\boxempty^\nu (A \at 1 \rightarrow C \at 1)) \at 1) \at 1 \\
  M^{\times} &= \lambda p.\, \llet_{1} (x, y) = p \lin
    \bbox_\nu\, (\unbox_\nu\, x, \unbox_\nu\, y) \\
    &\phantom{=}\; : ((\boxempty^\nu A \times \boxempty^\nu B) \at 1 \rightarrow \boxempty^\nu (A \times B) \at 1) \at 1 \\
  M^{+} &= \lambda s.\, \case_{1}\, s\,
    \{x.\bbox_\nu\, (\inl\, (\unbox_\nu\, x));\; \\
    &\phantom{= \lambda s.\, \case_{1}\, s\, \{}
    y.\bbox_\nu\, (\inr\, (\unbox_\nu\, y))\} \\
    &\phantom{=}\; : ((\boxempty^\nu A + \boxempty^\nu B) \at 1 \rightarrow \boxempty^\nu (A + B) \at 1) \at 1
\end{align*}

\subsection{Converting Graded Modal Type Theory to Mode Qualifiers}

\paragraph{Interpretation of types}
%
\begin{align*}
  \sem{\unit} &= \unit \\
  \sem{A + B} &= \sem{A} + \sem{B} \\
  \sem{A \times B} &= \sem{A} \times \sem{B} \\
  \sem{\boxempty^{q} A} &= \boxempty^{q}\, \sem{A} \\
  \sem{\,^q A \rightarrow B} &= \sem{A} \at q \rightarrow \sem{B} \at 1
\end{align*}

\paragraph{Contexts}
%
\begin{align*}
  \sem{\cdot} &= \cdot \\
  \sem{\Gamma, x :^0 A} &= \sem{\Gamma} \\
  \sem{\Gamma, x :^q A} &= \sem{\Gamma},\, x : \sem{A} \at q
\end{align*}

\paragraph{Translation of values}
%
\begin{align*}
  \sem{x} &= x \\
  \sem{()} &= () \\
  \sem{\lambda x.\, e} &= \lambda x.\, \sem{e} \\
  \sem{\bbox_q\, v} &= \bbox_q\; \sem{v} \\
  \sem{(v_1, v_2)} &= (\sem{v_1}, \sem{v_2}) \\
  \sem{\inl\, v} &= \inl\; \sem{v} \\
  \sem{\inr\, v} &= \inr\; \sem{v}
\end{align*}

\paragraph{Translation of expressions}
%
\begin{align*}
  \sem{v} &= \sem{v} \\
  \sem{e\, v} &= \sem{e}\; \sem{v} \\
  \sem{\llet x \leftarrow e_1 \lin e_2}
    &= \llet x = \sem{e_1} \lin \sem{e_2} \\
  \sem{\llet_{q_1} \bbox_{q_2} x = v \lin e}
    &= \llet x = \unbox_{q_2}\; \sem{v} \lin \sem{e} \\
  \sem{\llet_q (x, y) = v \lin e}
    &= \llet_q (x, y) = \sem{v} \lin \sem{e} \\
  \sem{\case_q\, v\, \{x.e_1;\; y.e_2\}}
    &= \case_q\; \sem{v}\; \{x.\sem{e_1};\; y.\sem{e_2}\}
\end{align*}

\begin{conjecture}[Graded Type Theory to Mode Qualifiers]
  If\/ $\Gamma \tee v : A$, then $\sem{\Gamma} \tee \sem{v} : \sem{A} \at 1$.
  If\/ $\Gamma \tee e : A$, then $\sem{\Gamma} \tee \sem{e} : \sem{A} \at 1$.
\end{conjecture}

\begin{acks}

The internship was supervised by Anton Lorenzen and Sam Lindley.

\end{acks}

\bibliographystyle{ACM-Reference-Format}
\bibliography{references.bib}

\end{document}
\endinput
